\documentclass[../NDCNNAI_main.tex]{subfiles}
\graphicspath{{\subfix{../fig/}}}

\begin{document}
\subsection{Метрические пространства, нормы, банаховы пространства}

\begin{definition}[Метрическое пространство]
    Пусть \( X \) --- произвольное множество. Функция \( d: X \times X \to \mathbb{R} \) называется \textbf{метрикой} на \( X \), если для любых элементов \( x, y, z \in X \) выполняются условия:
    \begin{enumerate}
        \item \( d(x, y) \ge 0 \) (неотрицательность),
        \item \( d(x, y) = 0 \iff x = y \) (тождественность неразличимых),
        \item \( d(x, y) = d(y, x) \) (симметричность),
        \item \( d(x, y) \le d(x, z) + d(z, y) \) (неравенство треугольника).
    \end{enumerate}
    Пара \( (X, d) \) называется \textbf{метрическим пространством}.
\end{definition}

\begin{definition}[Нормированное пространство]
    Пусть \( X \) --- векторное пространство над \( \mathbb{R} \). Функция \( \|\cdot\|: X \to \mathbb{R} \) называется \textbf{нормой} на \( X \), если для любых \( x, y \in X \) и \( \lambda \in \mathbb{R} \) выполняются:
    \begin{enumerate}
        \item \( \|x\| \ge 0 \), причём \( \|x\| = 0 \iff x = 0 \) (положительная определённость),
        \item \( \|\lambda x\| = |\lambda| \|x\| \) (однородность),
        \item \( \|x + y\| \le \|x\| + \|y\| \) (неравенство треугольника).
    \end{enumerate}
    Пара \( (X, \|\cdot\|) \) называется \textbf{нормированным пространством}. Норма индуцирует метрику: \( d(x, y) = \|x - y\| \).
\end{definition}

\begin{definition}[Банахово пространство]
    Нормированное пространство \( (X, \|\cdot\|) \) называется \textbf{банаховым}, если оно полно относительно индуцированной метрики, т.е. любая последовательность Коши в \( X \) сходится к элементу из \( X \).
\end{definition}
Примеры нормированных пространств:
    \begin{itemize}
        \item \( C(K) \) --- пространство непрерывных функций на компакте \( K \subset \mathbb{R}^n \) с равномерной нормой
              \[ \|f\|_\infty = \sup_{x \in K} |f(x)|. \]
        \item \( L^p(\mathbb{R}^n) \) --- пространство измеримых функций с конечной нормой
              \[ \|f\|_p = \left( \int_{\mathbb{R}^n} |f(x)|^p \, dx \right)^{1/p}, \quad 1 \le p < \infty. \]
        \item \( L^\infty(\mathbb{R}^n) \) --- пространство существенно ограниченных функций с нормой
              \[ \|f\|_\infty = \operatorname{ess\,sup}_{x \in \mathbb{R}^n} |f(x)| = \inf\{ C \ge 0 : |f(x)| \le C \text{ почти всюду} \}. \]
        \item \( L^1_{\text{loc}}(\mathbb{R}^n) \) --- пространство локально интегрируемых функций, для которых
              \[ \int_K |f(x)| \, dx < \infty \quad \text{для любого компакта } K \subset \mathbb{R}^n. \]
              Это пространство не является нормированным, но метризуемо (например, метрикой \( d(f,g) = \sum_{k=1}^\infty 2^{-k} \frac{\|f-g\|_{L^1(B_k)}}{1+\|f-g\|_{L^1(B_k)}} \), где \( B_k \) --- шар радиуса \( k \)).
    \end{itemize}


\subsection{Зачем функциональные пространства в теории ИНС и УА}

Нейронные сети приближают \textbf{функции} и \textbf{операторы} между функциональными пространствами. Для точной формулировки теорем универсальной аппроксимации (УА), устойчивости и оценки погрешности необходимы:

\begin{itemize}
    \item Строго определённые пространства (метрические, нормированные, банаховы),
    \item Нормы и двойственные пространства,
    \item Инструменты функционального анализа (теоремы Лебега, Хана–Банаха, \\
    Рисса–Маркова–Какутани, преобразование Фурье).
\end{itemize}
Напомним определения и теоремы из курсов математического и функционального анализа:

\begin{definition}[Полунорма]
Функция \( p\colon X \to \mathbb{R} \) на векторном пространстве \( X \) называется \textbf{полунормой}, если для всех \( x, y \in X \) и \( \lambda \in \mathbb{R} \):
\begin{enumerate}
    \item \( p(x) \ge 0 \) (неотрицательность),
    \item \( p(\lambda x) = |\lambda| p(x) \) (однородность),
    \item \( p(x + y) \le p(x) + p(y) \) (неравенство треугольника).
\end{enumerate}
Отличие от нормы: из \( p(x) = 0 \) не следует \( x = 0 \).
\end{definition}

\begin{definition}[Борелевская мера]
Пусть \( X \) --- топологическое пространство. Мера \( \mu \), определённая на \textbf{борелевской \( \sigma \)-алгебре} (наименьшей \( \sigma \)-алгебре, содержащей все открытые множества), называется \textbf{борелевской мерой}.
\end{definition}

\begin{definition}[Регулярная борелевская мера]
Борелевская мера \( \mu \) на локально компактном хаусдорфовом пространстве \( X \) называется \textbf{регулярной}, если имеются:
\begin{enumerate}
    \item \textbf{Внутренняя регулярность:} Для любого борелевского множества \( E \subset X \),
    \[
    \mu(E) = \sup\{\mu(K)\colon K \subset E, \ K \text{ компакт}\}.
    \]
    \item \textbf{Внешняя регулярность:} Для любого борелевского множества \( E \subset X \),
    \[
    \mu(E) = \inf\{\mu(U)\colon E \subset U, \ U \text{ открыто}\}.
    \]
\end{enumerate}
\end{definition}

\begin{definition}[Топологическое двойственное пространство]
Пусть \( X \) --- топологическое векторное пространство. \textbf{Двойственное (сопряжённое) пространство} \( X^* \) --- это множество всех непрерывных линейных функционалов \( f\colon X \to \mathbb{R} \). Если \( X \) --- нормированное пространство, то \( X^* \) снабжается нормой
\[
\|f\|_{X^*} = \sup_{\|x\|_X \le 1} |f(x)|.
\]
\end{definition}
\begin{definition}[Преобразование Фурье в $L^1(\mathbb{R}^n)$]
Для функции $f \in L^1(\mathbb{R}^n)$ её \textbf{преобразование Фурье} определяется как
\[
\hat{f}(\xi) = \mathcal{F}[f](\xi) = \int_{\mathbb{R}^n} f(x) e^{-2\pi i \xi \cdot x} \, dx, \quad \xi \in \mathbb{R}^n.
\]
При этом $\hat{f} \in L^\infty(\mathbb{R}^n)$ и $\|\hat{f}\|_\infty \le \|f\|_1$.
\end{definition}

\begin{definition}[Преобразование Фурье в $L^2(\mathbb{R}^n)$]
Для $f \in L^2(\mathbb{R}^n)$ преобразование Фурье можно определить как предел в $L^2$-норме:
\[
\hat{f}(\xi) = \lim_{R \to \infty} \int_{|x| \le R} f(x) e^{-2\pi i \xi \cdot x} \, dx.
\]
Оператор $\mathcal{F}\colon L^2(\mathbb{R}^n) \to L^2(\mathbb{R}^n)$ является унитарным (сохраняет $L^2$-норму).
\end{definition}

\begin{definition}[Преобразование Фурье меры]
Для конечной комплекснозначной борелевской меры $\mu$ на $\mathbb{R}^n$ её \textbf{преобразование Фурье} (характеристическая функция) определяется как
\[
\hat{\mu}(\xi) = \int_{\mathbb{R}^n} e^{-2\pi i \xi \cdot x} \, d\mu(x).
\]
\end{definition}

\begin{theorem}[Обращение преобразования Фурье]
Если $f, \hat{f} \in L^1(\mathbb{R}^n)$, то справедлива формула обращения:
\[
f(x) = \int_{\mathbb{R}^n} \hat{f}(\xi) e^{2\pi i \xi \cdot x} \, d\xi \quad \text{п.в.}
\]
\end{theorem}

\begin{theorem}[Свойства преобразования Фурье]
Пусть $f, g \in L^1(\mathbb{R}^n)$. Тогда:
\begin{enumerate}
    \item \textbf{Линейность:} $\mathcal{F}[\alpha f + \beta g] = \alpha \hat{f} + \beta \hat{g}$.
    \item \textbf{Сдвиг:} если $g(x) = f(x - a)$, то $\hat{g}(\xi) = e^{-2\pi i a \cdot \xi} \hat{f}(\xi)$.
    \item \textbf{Модуляция:} если $g(x) = e^{2\pi i a \cdot x} f(x)$, то $\hat{g}(\xi) = \hat{f}(\xi - a)$.
    \item \textbf{Дифференцирование:} если $\partial_{x_j} f \in L^1$, то $\widehat{\partial_{x_j} f}(\xi) = 2\pi i \xi_j \hat{f}(\xi)$.
\end{enumerate}
\end{theorem}
\begin{theorem}[Хана–Банаха, аналитическая форма]
Пусть \( X \) --- вещественное векторное пространство, \( p \) --- полунорма на \( X \), \( L \) --- линейное подпространство, \( f\colon L \to \mathbb{R} \) --- линейный функционал, удовлетворяющий \( |f(x)| \le p(x) \) для всех \( x \in L \). Тогда существует линейный функционал \( F\colon X \to \mathbb{R} \), продолжающий \( f \) с сохранением подчинённости полунорме: \( |F(x)| \le p(x) \) для всех \( x \in X \).
\end{theorem}

\begin{theorem}[Рисса–Маркова–Какутани]
Пусть \( X \) --- локально компактное хаусдорфово пространство. Тогда всякий положительный линейный функционал \( \Lambda\colon C_c(X) \to \mathbb{R} \) (где \( C_c(X) \) --- непрерывные функции с компактным носителем) представляется в виде
\[
\Lambda(f) = \int_X f \, d\mu
\]
для некоторой единственной регулярной борелевской меры \( \mu \) на \( X \).
\end{theorem}

\begin{theorem}[Лебега, о мажорированной сходимости]
Пусть \( \{f_n\} \) --- последовательность измеримых функций на \( \Omega \), сходящаяся почти всюду к \( f \). Если существует интегрируемая функция \( g \in L^1(\Omega) \) такая, что \( |f_n(x)| \le g(x) \) для всех \( n \) и почти всех \( x \), то \( f \in L^1(\Omega) \) и
\[
\lim_{n\to\infty} \int_\Omega f_n \, d\mu = \int_\Omega f \, d\mu.
\]
\end{theorem}
Выбор пространства (например, \( C(K) \), \( L^p(\mu) \), \( L^1_{\text{loc}} \)) определяет:
\begin{itemize}
    \item Тип сходимости (равномерная, в среднем (сходимость по норме \( L^p \)), поточечная),
    \item Условия на активационную функцию,
    \item Наличие порогового члена (сдвига \( \theta \)) для полноты аппроксимации.
\end{itemize}

\begin{definition}[Ridge-функция]
    Функция вида \( f(x) = \sigma(w^T x + \theta) \), где \( \sigma\colon \mathbb{R} \to \mathbb{R} \), \( w \in \mathbb{R}^n \), \( \theta \in \mathbb{R} \), называется \textbf{ridge-функцией}.
\end{definition}

\begin{definition}[Сигмоидальная функция]
    Функция \( \sigma\colon \mathbb{R} \to \mathbb{R} \) называется \textbf{сигмоидальной}, если
    \[
    \lim_{t \to +\infty} \sigma(t) = 1 \quad \text{и} \quad \lim_{t \to -\infty} \sigma(t) = 0.
    \]
    Типичный пример: логистическая функция \( \sigma(t) = \frac{1}{1+e^{-t}} \).
\end{definition}

\begin{definition}[Дискриминаторная функция]
    Ограниченная измеримая функция \( \sigma \) называется \textbf{дискриминаторной}, если для любой конечной знакопеременной меры (заряда) \( \mu \) на компакте \( K \subset \mathbb{R}^n \) условие
    \[
    \int_K \sigma(w^T x + \theta) \, d\mu(x) = 0 \quad \forall w \in \mathbb{R}^n, \theta \in \mathbb{R}
    \]
    влечёт \( \mu = 0 \).
\end{definition}
Приведём ряд основных результатов:

\begin{theorem}[Цыбенко]
    Пусть \( \sigma \) --- непрерывная сигмоидальная функция. Тогда для любой непрерывной функции \( f \in C(K) \) на компакте \( K \subset \mathbb{R}^n \) и любого \( \varepsilon > 0 \) существует функция вида
    \[
    G(x) = \sum_{j=1}^{N} \alpha_j \sigma(w_j^T x + \theta_j),
    \]
    такая что
    \[
    \sup_{x \in K} |G(x) - f(x)| < \varepsilon.
    \]
\end{theorem}

\begin{remark}
    Функции вида \( G(x) \) соответствуют выходу однослойной нейронной сети с \( N \) скрытыми нейронами.
\end{remark}
\begin{theorem}[Обобщение на \( L^p \)]
    Если \( \sigma \) ограничена и непостоянна, то для любой конечной меры \( \mu \) на \( \mathbb{R}^n \) суммы вида
    \[
    G(x) = \sum_{j=1}^{N} \alpha_j \sigma(w_j^T x + \theta_j)
    \]
    плотны в \( L^p(\mu) \), \( 1 \le p < \infty \).
\end{theorem}

\begin{theorem}[Обобщение на \( C(K) \)]
    Если \( \sigma \) непрерывна, ограничена и непостоянна, то для любого компакта \( K \subset \mathbb{R}^n \) суммы вида \( G(x) \) плотны в \( C(K) \) в равномерной норме.
\end{theorem}

\begin{remark}
    Таким образом, для универсальной аппроксимации в \( L^p \) или \( C(K) \) достаточно, чтобы \( \sigma \) была \textbf{ограниченной и непостоянной} (не обязательно сигмоидальной).
\end{remark}

\begin{theorem}[Необходимое и достаточное условие]
    Пусть \( \sigma \) --- локально ограниченная функция, множество точек разрыва которой имеет нулевую меру Лебега. Тогда
    \[
    \Sigma_n = \operatorname{Lin}\{\sigma(w^T x + \theta)\colon w \in \mathbb{R}^n, \theta \in \mathbb{R}\}
    \]
    плотно в \( C(\mathbb{R}^n) \) \textbf{тогда и только тогда}, когда \( \sigma \) \textbf{не является полиномом} (почти всюду).
\end{theorem}

\begin{remark}
    Это самый общий результат: активационная функция может быть разрывной, неограниченной, но \textbf{не должна быть полиномом}. При этом наличие порога \( \theta \) существенно.
\end{remark}
\begin{itemize}
    \item Без порога \( \theta \) утверждение последней теоремы неверно. Пример: \( \sigma(x) = \sin(x) \) без порога не может аппроксимировать чётные функции.
    \item Порог обеспечивает сдвиг активационной функции, что необходимо для полноты класса ridge-функций \( \sigma(w^T x + \theta) \).
\end{itemize}
Таким образом, получаем:
\begin{itemize}
    \item Функциональные пространства \( C(K) \), \( L^p \), \( L^\infty \) задают естественную метрику для анализа аппроксимационных свойств нейронных сетей.
    \item Теорема Цыбенко показывает, что однослойные сети с сигмоидальной активацией могут приближать любые непрерывные функции на компактах.
    \item Последующие обобщения показывают, что достаточно ограниченности и непостоянства, а в общем случае --- неполнотности полиномами.
    \item \textbf{Ограничение:} в \( L^\infty \) универсальная аппроксимация невозможна из-за несепарабельности и природы разрывных функций.
\end{itemize}

\end{document}